% GENERATED FILE. Edit canonical lessons, then run npm run generate:books.
% canonical-source-hash: fnv1a64:c94fefa1f96aaffe
% canonical-lessons: RU-C326-predok, RU-C326-potomok, RU-C326-letopis, RU-C326-knyaz, RU-C326-osnovyvat, RU-R326-review-319-326

\chapter{Ancestors and chronicles}
\label{ch:ru-ancestors-and-chronicles}
\hlchaptermodality{326}

\begin{chapteropening}
\textbf{By the end of this chapter:} \emph{I can talk about an ancestor, a descendant, a chronicle and a prince of old Rus, and say that a town was founded.}

Complete the last lesson of chapter 326: I can talk about an ancestor, a descendant, a chronicle and a prince of old Rus, and say that a town was founded.
\end{chapteropening}

\section[\texorpdfstring{\ru{предок} (prédok)}{предок (prédok)}]{\ru{предок} (prédok) — an ancestor}

\label{lesson:RU-C326-predok}



\begin{quote}
Before the new one: say the Russian for to conquer, then the Russian for to rule.
\end{quote}

\subsection*{You'll want to know: \ru{предок}}
\textbf{\ru{предок}} (\emph{prédok}) — \textquotedblleft{}an ancestor\textquotedblright{}.

\begin{quote}
\ru{Наши} \ru{предки} \ru{жили} \ru{в} \ru{этой} \ru{деревне}.
\end{quote}

\begin{quote}
\emph{Náshi prédki zhíli v étoy derévne.}
\end{quote}

\begin{quote}
Our ancestors lived in this village.
\end{quote}

Genitive \emph{\ru{предка}}. In slang \emph{\ru{предки}} are one's parents.

\begin{cousinweb}
\emph{\ru{предок}} is a Church Slavonic word, \emph{prědŭkŭ} (EtymDB), built on \emph{\ru{пред}-} \textquotedblleft{}before\textquotedblright{}, the church shape of native \emph{\ru{перед}}: the one who went before.
\end{cousinweb}

\begin{grammarlens}[title={What you've built}]
One more word for telling, explaining and describing at B1.
\end{grammarlens}

\subsection*{Your turn}
\begin{itemize}
\raggedright
  \item \emph{Say it:} \emph{prédok}
  \item \emph{Say it:} \emph{prédok}, once more
  \item \emph{Say it:} a sentence of your own with \emph{prédok}
  \item \emph{Recall:} say the Russian for to conquer, then the Russian for to rule, then say \emph{prédok} again
\end{itemize}

\begin{tcolorbox}[breakable,colback=teal!4,colframe=teal!35!black,title={Before you move on}]
What does \ru{предок} mean? (\textquotedblleft{}An ancestor\textquotedblright{}.) Say it once more.
\end{tcolorbox}
\section[\texorpdfstring{\ru{потомок} (potómok)}{потомок (potómok)}]{\ru{потомок} (potómok) — a descendant}

\label{lesson:RU-C326-potomok}



\begin{quote}
Before the new one: say the Russian for to rule, then the Russian for an ancestor.
\end{quote}

\subsection*{You'll want to know: \ru{потомок}}
\textbf{\ru{потомок}} (\emph{potómok}) — \textquotedblleft{}a descendant\textquotedblright{}.

\begin{quote}
\ru{Он} \ru{потомок} \ru{известного} \ru{художника}.
\end{quote}

\begin{quote}
\emph{On potómok izvéstnovo khudózhnika.}
\end{quote}

\begin{quote}
He is a descendant of a famous painter.
\end{quote}

Genitive \emph{\ru{потомка}}.

\begin{cousinweb}
\emph{\ru{потомок}} is usually traced to \emph{\ru{потом}} \textquotedblleft{}afterwards\textquotedblright{}, from \emph{\ru{по}} and \emph{\ru{том}}, a form of \emph{\ru{тот}} \textquotedblleft{}that\textquotedblright{}: the one who comes after. EtymDB lists no parent.
\end{cousinweb}

\begin{grammarlens}[title={What you've built}]
One more word for telling, explaining and describing at B1.
\end{grammarlens}

\subsection*{Your turn}
\begin{itemize}
\raggedright
  \item \emph{Say it:} \emph{potómok}
  \item \emph{Say it:} \emph{potómok}, once more
  \item \emph{Say it:} a sentence of your own with \emph{potómok}
  \item \emph{Recall:} say the Russian for to rule, then the Russian for an ancestor, then say \emph{potómok} again
\end{itemize}

\begin{tcolorbox}[breakable,colback=teal!4,colframe=teal!35!black,title={Before you move on}]
What does \ru{потомок} mean? (\textquotedblleft{}A descendant\textquotedblright{}.) Say it once more.
\end{tcolorbox}
\section[\texorpdfstring{\ru{летопись} (létopis)}{летопись (létopis)}]{\ru{летопись} (létopis) — a chronicle, annals}

\label{lesson:RU-C326-letopis}



\begin{quote}
Before the new one: say the Russian for an ancestor, then the Russian for a descendant.
\end{quote}

\subsection*{You'll want to know: \ru{летопись}}
\textbf{\ru{летопись}} (\emph{létopis}) — \textquotedblleft{}a chronicle, annals\textquotedblright{}.

\begin{quote}
\ru{Город} \ru{впервые} \ru{упоминается} \ru{в} \ru{летописи}.
\end{quote}

\begin{quote}
\emph{Górod vpervýye upomináyetsya v létopisi.}
\end{quote}

\begin{quote}
The town is first mentioned in a chronicle.
\end{quote}

Feminine. The old chronicles of Rus were written year by year.

\begin{cousinweb}
\emph{\ru{летопись}} joins \emph{\ru{лето}} \textquotedblleft{}summer, year\textquotedblright{}, Proto-Slavic \emph{lěto} \textquotedblleft{}year\textquotedblright{}, and \emph{\ru{писать}} \textquotedblleft{}to write\textquotedblright{}, Proto-Slavic \emph{pisati} (EtymDB): a record kept year by year.
\end{cousinweb}

\begin{grammarlens}[title={What you've built}]
One more word for telling, explaining and describing at B1.
\end{grammarlens}

\subsection*{Your turn}
\begin{itemize}
\raggedright
  \item \emph{Say it:} \emph{létopis}
  \item \emph{Say it:} \emph{létopis}, once more
  \item \emph{Say it:} a sentence of your own with \emph{létopis}
  \item \emph{Recall:} say the Russian for an ancestor, then the Russian for a descendant, then say \emph{létopis} again
\end{itemize}

\begin{tcolorbox}[breakable,colback=teal!4,colframe=teal!35!black,title={Before you move on}]
What does \ru{летопись} mean? (\textquotedblleft{}A chronicle, annals\textquotedblright{}.) Say it once more.
\end{tcolorbox}
\section[\texorpdfstring{\ru{князь} (knyaz)}{князь (knyaz)}]{\ru{князь} (knyaz) — a prince of old Rus, a ruling prince}

\label{lesson:RU-C326-knyaz}



\begin{quote}
Before the new one: say the Russian for a descendant, then the Russian for a chronicle.
\end{quote}

\subsection*{You'll want to know: \ru{князь}}
\textbf{\ru{князь}} (\emph{knyaz}) — \textquotedblleft{}a prince of old Rus, a ruling prince\textquotedblright{}.

\begin{quote}
\ru{Князь} \ru{жил} \ru{в} \ru{большом} \ru{доме} \ru{на} \ru{холме}.
\end{quote}

\begin{quote}
\emph{Knyaz zhil v bolshóm dóme na kholmé.}
\end{quote}

\begin{quote}
The prince lived in a big house on a hill.
\end{quote}

Plural \emph{\ru{князья}}. \emph{\ru{Княгиня}} — a princess, a prince's wife.

\begin{cousinweb}
\emph{\ru{князь}} is inherited from Proto-Slavic \emph{kŭnędzĭ}, an early borrowing from Proto-Germanic \emph{kuningaz} \textquotedblleft{}king\textquotedblright{} (EtymDB); English \emph{king} is a cousin.
\end{cousinweb}

\begin{grammarlens}[title={What you've built}]
One more word for telling, explaining and describing at B1.
\end{grammarlens}

\subsection*{Your turn}
\begin{itemize}
\raggedright
  \item \emph{Say it:} \emph{knyaz}
  \item \emph{Say it:} \emph{knyaz}, once more
  \item \emph{Say it:} a sentence of your own with \emph{knyaz}
  \item \emph{Recall:} say the Russian for a descendant, then the Russian for a chronicle, then say \emph{knyaz} again
\end{itemize}

\begin{tcolorbox}[breakable,colback=teal!4,colframe=teal!35!black,title={Before you move on}]
What does \ru{князь} mean? (\textquotedblleft{}A prince of old Rus, a ruling prince\textquotedblright{}.) Say it once more.
\end{tcolorbox}
\section[\texorpdfstring{\ru{основывать} (osnóvyvat)}{основывать (osnóvyvat)}]{\ru{основывать} (osnóvyvat) — to found}

\label{lesson:RU-C326-osnovyvat}



\begin{quote}
Before the new one: say the Russian for a chronicle, then the Russian for a prince of old Rus.
\end{quote}

\subsection*{You'll want to know: \ru{основывать}}
\textbf{\ru{основывать}} (\emph{osnóvyvat}) — \textquotedblleft{}to found\textquotedblright{}.

\begin{quote}
\ru{Город} \ru{основали} \ru{в} \ru{двенадцатом} \ru{веке}.
\end{quote}

\begin{quote}
\emph{Górod osnováli v dvenádtsatom véke.}
\end{quote}

\begin{quote}
The town was founded in the twelfth century.
\end{quote}

Imperfective; perfective \emph{\ru{основать}} (\emph{osnovát}). \emph{\ru{Основатель}} — a founder.

\begin{cousinweb}
\emph{\ru{основывать}} is built on \emph{\ru{о}-} and \emph{\ru{основа}} \textquotedblleft{}a base\textquotedblright{}, inherited from Proto-Slavic \emph{osnova} (EtymDB), first the warp threads on a loom, a cousin of \emph{\ru{сновать}} \textquotedblleft{}to warp, to dart to and fro\textquotedblright{}.
\end{cousinweb}

\begin{grammarlens}[title={What you've built}]
One more word for telling, explaining and describing at B1.
\end{grammarlens}

\subsection*{Your turn}
\begin{itemize}
\raggedright
  \item \emph{Say it:} \emph{osnóvyvat}
  \item \emph{Say it:} \emph{osnóvyvat}, once more
  \item \emph{Say it:} a sentence of your own with \emph{osnóvyvat}
  \item \emph{Recall:} say the Russian for a chronicle, then the Russian for a prince of old Rus, then say \emph{osnóvyvat} again
\end{itemize}

\begin{tcolorbox}[breakable,colback=teal!4,colframe=teal!35!black,title={Before you move on}]
What does \ru{основывать} mean? (\textquotedblleft{}To found\textquotedblright{}.) Say it once more.
\end{tcolorbox}
\section[(dialogue)]{Words from chapters 319–326 — a review}

\label{lesson:RU-R326-review-319-326}



\begin{quote}
Say the Russian for a legend and a myth.
\end{quote}

\subsection*{Your turn}
Say these aloud:

\begin{itemize}
\raggedright
  \item a legend, a myth, a hero, a storyteller, to make up
  \item a wizard, a witch, a monster, a dragon, to turn into
  \item a king, a queen, a princess, a prince, a kingdom
  \item a knight, armour, a sword, a shield, a throne
  \item a war, a battle, a pitched battle, to fight, an enemy
  \item a victory, to win, to wage war, to die, to free
  \item an empire, a revolution, an uprising, to conquer, to rule
  \item an ancestor, a descendant, a chronicle, a prince of old Rus, to found
\end{itemize}

\begin{tcolorbox}[breakable,colback=teal!4,colframe=teal!35!black,title={Before you move on}]
A legend? (\textbf{\ru{легенда}}.) A myth? (\textbf{\ru{миф}}.) A hero? (\textbf{\ru{герой}}.) A storyteller? (\textbf{\ru{рассказчик}}.) To make up? (\textbf{\ru{сочинять}}.) A wizard? (\textbf{\ru{волшебник}}.) A witch? (\textbf{\ru{ведьма}}.) A monster? (\textbf{\ru{чудовище}}.) A dragon? (\textbf{\ru{дракон}}.) To turn into? (\textbf{\ru{превращаться}}.) A king? (\textbf{\ru{король}}.) A queen? (\textbf{\ru{королева}}.) A princess? (\textbf{\ru{принцесса}}.) A prince? (\textbf{\ru{принц}}.) A kingdom? (\textbf{\ru{королевство}}.) A knight? (\textbf{\ru{рыцарь}}.) Armour? (\textbf{\ru{доспехи}}.) A sword? (\textbf{\ru{меч}}.) A shield? (\textbf{\ru{щит}}.) A throne? (\textbf{\ru{трон}}.) A war? (\textbf{\ru{война}}.) A battle? (\textbf{\ru{битва}}.) A pitched battle? (\textbf{\ru{сражение}}.) To fight? (\textbf{\ru{сражаться}}.) An enemy? (\textbf{\ru{враг}}.) A victory? (\textbf{\ru{победа}}.) To win? (\textbf{\ru{побеждать}}.) To wage war? (\textbf{\ru{воевать}}.) To die? (\textbf{\ru{погибать}}.) To free? (\textbf{\ru{освобождать}}.) An empire? (\textbf{\ru{империя}}.) A revolution? (\textbf{\ru{революция}}.) An uprising? (\textbf{\ru{восстание}}.) To conquer? (\textbf{\ru{завоёвывать}}.) To rule? (\textbf{\ru{править}}.) An ancestor? (\textbf{\ru{предок}}.) A descendant? (\textbf{\ru{потомок}}.) A chronicle? (\textbf{\ru{летопись}}.) A prince of old Rus? (\textbf{\ru{князь}}.) To found? (\textbf{\ru{основывать}}.)
\end{tcolorbox}
